\Article{38}
Wir gehen nun zur Kongruenz von Zahlen nach komplexen Moduln über. Zu Beginn dieser Untersuchung ist jedoch anzugeben, wie sich der Bereich der komplexen Größen anschaulich darstellen lässt.

Jede reelle Größe lässt sich durch eine Strecke auf einer nach beiden Seiten unendlichen Geraden ausdrücken, die von einem beliebigen Anfangspunkt aus genommen und nach einer beliebigen, als Einheit gewählten Strecke gemessen wird. Sie lässt sich daher durch den anderen Endpunkt darstellen, wobei die Punkte auf der einen Seite des Anfangspunktes positive, auf der anderen negative Größen darstellen. Ebenso lässt sich jede komplexe Größe durch einen Punkt in einer unendlichen Ebene darstellen, in der eine bestimmte Gerade den reellen Größen zugeordnet ist: die komplexe Größe $x+iy$ durch den Punkt mit der Abszisse $=x$ und der Ordinate $=y$, wobei diese auf der einen Seite der Abszissenlinie positiv, auf der anderen negativ genommen wird. So kann man sagen, dass jede komplexe Größe den Unterschied zwischen der Lage des ihr zugeordneten Punktes und der Lage des Anfangspunktes misst. Die positive Einheit bezeichnet dabei eine beliebige festgelegte Verschiebung in einer beliebigen festgelegten Richtung, die negative Einheit eine gleich große Verschiebung in der entgegengesetzten Richtung und die imaginären Einheiten gleich große Verschiebungen in den beiden senkrechten seitlichen Richtungen.

Auf diese Weise wird die Metaphysik der Größen, die wir imaginär nennen, erheblich erhellt. Wird der Anfangspunkt mit $(0)$ bezeichnet und werden zwei komplexe Größen $m,m'$ den Punkten $M,M'$ zugeordnet, deren Lage bezüglich $(0)$ sie ausdrücken, so ist die Differenz $m-m'$ nichts anderes als die Lage des Punktes $M$ bezüglich des Punktes $M'$. Stellt dagegen das Produkt $mm'$ die Lage des Punktes $N$ bezüglich $(0)$ dar, so sieht man leicht, dass diese Lage durch die Lage des Punktes $M$ bezüglich $(0)$ ebenso bestimmt wird wie die Lage des Punktes $M'$ durch die Lage des der positiven Einheit entsprechenden Punktes. Man kann daher passend sagen, dass die Lagen der den komplexen Größen $mm'$
\SourceLeaf{120}%
$m,m',1$ entsprechenden Punkte eine \emph{Proportion} bilden. Eine ausführlichere Behandlung behalten wir uns jedoch für eine andere Gelegenheit vor. Die Schwierigkeiten, in die man die Theorie der imaginären Größen verstrickt glaubt, sind großenteils aus ungeeigneten Benennungen entstanden (haben doch manche sogar den missklingenden Namen unmöglicher Größen gebraucht). Wären wir von den Begriffen ausgegangen, die zweidimensionale Mannigfaltigkeiten darbieten, wie sie in größter Reinheit in unseren Raumanschauungen erscheinen, und hätten wir die positiven Größen direkte, die negativen inverse und die imaginären seitliche Größen genannt, so wäre an die Stelle von Verwicklungen Einfachheit und an die Stelle von Dunkelheit Klarheit getreten.

\Article{39}
Das im vorigen Artikel Dargelegte bezieht sich auf stetige komplexe Größen. In der Arithmetik, die sich nur mit ganzen Zahlen befasst, wird das Bild der komplexen Zahlen ein System gleich weit voneinander entfernter Punkte sein, die auf gleich weit voneinander entfernten Geraden so angeordnet sind, dass sie die unendliche Ebene in unendlich viele gleiche Quadrate zerlegen. Alle durch eine gegebene komplexe Zahl $a+bi=m$ teilbaren Zahlen bilden ebenfalls unendlich viele Quadrate, mit den Seiten $=\surd(aa+bb)$ oder den Flächeninhalten $=aa+bb$. Die letzteren Quadrate sind gegen die ersteren geneigt, sofern keine der beiden Zahlen $a,b$ $=0$ ist. Jeder durch den Modul $m$ nicht teilbaren Zahl entspricht ein Punkt entweder innerhalb eines solchen Quadrates oder auf einer zwei Quadraten gemeinsamen Seite; der letztere Fall kann jedoch nur eintreten, wenn $a,b$ einen gemeinsamen Teiler haben. Ferner ist klar, dass nach dem Modul $m$ kongruente Zahlen in ihren Quadraten entsprechende Lagen einnehmen. Sammeln wir daher alle Zahlen innerhalb eines bestimmten Quadrates sowie alle, die etwa auf zwei einander nicht gegenüberliegenden Seiten liegen, und fügen schließlich eine durch $m$ teilbare Zahl hinzu, so erhalten wir leicht ein vollständiges System inkongruenter Reste nach dem Modul $m$: Jede ganze Zahl muss einem von ihnen und nur einem einzigen kongruent sein. Es wäre auch nicht schwierig zu zeigen, dass die Anzahl dieser Reste der Norm des Moduls gleich ist, nämlich $=aa+bb$. Doch erscheint es ratsam, diesen sehr wichtigen Satz auf eine andere, rein arithmetische Weise zu beweisen.

\Article{40}
\textsc{Satz.}\quad\emph{Nach einem gegebenen komplexen Modul $m=a+bi$, dessen Norm $aa+bb=p$ ist und für den $a,b$ teilerfremd sind, ist jede komplexe ganze Zahl einem Rest aus der Folge $0,1,2,3\dotsfour p-1$ und nicht mehreren kongruent.}


\SourceLeaf{121}%
\emph{Beweis.} I. Es seien $\alpha,\beta$ ganze Zahlen mit $\alpha a+\beta b=1$. Dann ist
\[
i=\alpha b-\beta a+m(\beta+\alpha i)
\]
Für eine gegebene komplexe ganze Zahl $A+Bi$ haben wir daher
\[
A+Bi=A+(\alpha b-\beta a)B+m(\beta B+\alpha Bi)
\]
Bezeichnen wir mit $h$ den kleinsten positiven Rest der Zahl $A+(\alpha b-\beta a)B$ nach dem Modul $p$ und setzen
\[
A+(\alpha b-\beta a)B=h+kp=h+m(ak-bki)
\]
so ist
\[
A+Bi=h+m(\beta B+ak+(\alpha B-bk)i)
\]
oder
\[
A+Bi\equiv h(\mathrm{mod.}\ m).\qquad\text{Q. E. P.}
\]

II. Sind zwei reelle Zahlen $h,h'$ derselben komplexen Zahl nach dem Modul $m$ kongruent, so sind sie auch einander kongruent. Wir setzen daher $h-h'=m(c+di)$, woraus folgt:
\[
(h-h')(a-bi)=p(c+di)
\]
und somit
\[
(h-h')a=pc,\quad(h-h')b=-pd
\]
sowie wegen $a\alpha+b\beta=1$
\[
h-h'=p(c\alpha-d\beta),\ \text{d.\,h.}\ h\equiv h'(\mathrm{mod.}\ p)
\]
Sind also $h$ und $h'$ verschieden, so können sie nicht beide zugleich in der Sammlung der Zahlen $0,1,2,3\dotsfour p-1$ enthalten sein. Q. E. S.

\Article{41}
\textsc{Satz.}\quad\emph{Nach einem komplexen Modul $m=.\,a+bi$, dessen Norm $aa+bb=p$ ist und für den $a,b$ nicht teilerfremd sind, sondern den größten gemeinsamen Teiler $\lambda$ haben (den wir positiv genommen voraussetzen), ist jede komplexe Zahl einem Rest $x+yi$ kongruent, bei dem $x$ eine der Zahlen $0,1,2,3\dotsfour\frac p\lambda-1$ und $y$ eine der Zahlen $0,1,2,3\dotsfour\lambda-1$ ist, und zwar genau einem unter allen $p$ Resten dieser Form.}

\SourceLeaf{122}%
\emph{Beweis.} I. Wählen wir ganze Zahlen $\alpha,\beta$ mit $\alpha a+\beta b=\lambda$, so ist $\lambda i=\alpha b-\beta a+m(\beta+\alpha i)$. Es sei nun $A+Bi$ eine gegebene komplexe Zahl, $y$ der kleinste positive Rest von $B$ nach dem Modul $\lambda$ und $x$ der kleinste positive Rest von $A+(\alpha b-\beta a)\pdot\frac{B-y}{\lambda}$ nach dem Modul $\frac p\lambda$; wir setzen
\[
A+(\alpha b-\beta a)\pdot\frac{B-y}{\lambda}=x+\frac p\lambda\pdot k
\]
Dann ist
\[
\begin{aligned}
A+Bi-(x+yi)&=\frac p\lambda\pdot k+(B-y)i-(\alpha b-\beta a)\frac{B-y}{\lambda}\\
&=\frac p\lambda\pdot k+\frac{B-y}{\lambda}\pdot m(\beta+\alpha i)\\
&=\left(\frac a\lambda-\frac b\lambda\pdot i\right)km+\frac{B-y}{\lambda}(\beta+\alpha i)m
\end{aligned}
\]
also durch $m$ teilbar, oder gleichbedeutend $A+Bi\equiv x+yi(\mathrm{mod.}\ m)$ \mbox{Q. E. P.}

II. Wir nehmen an, dass nach dem Modul $m$ derselben komplexen Zahl zwei Zahlen $x+yi,x'+y'i$ kongruent sind, die daher auch einander nach dem Modul $m$ kongruent sind. Umso mehr sind sie nach dem Modul $\lambda$ kongruent, also gilt $y\equiv y'(\mathrm{mod.}\lambda)$. Werden daher beide $y,y'$ als unter den Zahlen $0,1,2,3\dotsfour\lambda-1$ enthalten vorausgesetzt, so muss $y=y'$ gelten. Damit gilt aber auch $x\equiv x'(\mathrm{mod.}\ m)$, das heißt, $x-x'$ ist durch $m$ teilbar und folglich die ganze Zahl $\frac{x-x'}\lambda$ durch $\frac a\lambda+\frac b\lambda\pdot i$, oder
\[
\frac{x-x'}\lambda\equiv0\left(\mathrm{mod.}\frac a\lambda+\frac b\lambda\pdot i\right)
\]
Da nun $\frac a\lambda,\frac b\lambda$ teilerfremd sind, folgt aus dem zweiten Teil des Satzes im vorigen Artikel, dass $\frac{x-x'}\lambda$ auch durch die Norm der Zahl $\frac a\lambda+\frac b\lambda\pdot i$, also durch die Zahl $\frac p{\lambda\lambda}$, teilbar ist, und somit $x-x'$ durch $\frac p\lambda$. Werden also auch beide $x,x'$ als in der Sammlung der Zahlen $0,1,2,3\dotsfour\frac p\lambda-1$ enthalten vorausgesetzt, so muss $x=x'$ gelten, oder die Reste $x+yi,x'+y'i$ sind identisch. Q. E. S.

Ferner ist offenbar auch der Fall hierher zu rechnen, in dem der Modul reell ist, also $b=0$ und folglich $\lambda=\pm a$, sowie der Fall, in dem der Modul rein imaginär ist, also $a=0$ und folglich $\lambda=\pm b$. In beiden Fällen gilt $\frac p\lambda=\lambda$.

\SourceLeaf{123}%
\Article{42}
Ordnet man somit alle nach einem gegebenen Modul einander kongruenten komplexen Zahlen derselben Klasse und inkongruente Zahlen verschiedenen Klassen zu, so gibt es insgesamt $p$ Klassen, die den ganzen Bereich der ganzen Zahlen erschöpfen, wobei $p$ die Norm des Moduls bezeichnet. Eine Sammlung ebenso vieler Zahlen, aus jeder Klasse eine genommen, stellt ein vollständiges System inkongruenter Reste dar, wie wir es in den Artikeln 40, 41 angegeben haben. In jenem System beruhte die Wahl der ihre Klassen gleichsam vertretenden Reste auf dem Prinzip, in jeder Klasse einen Rest $x+yi$ zu wählen, für den $y$ den kleinsten Wert hat, und unter allen mit demselben kleinsten Wert von $y$ denjenigen, für den $x$ den kleinsten Wert hat; dabei sind negative Werte sowohl für $x$ als auch für $y$ ausgeschlossen. Für andere Zwecke wird man jedoch andere Prinzipien verwenden. Besonders zu beachten ist die Wahl solcher Reste, die bei Division durch den Modul die einfachsten Quotienten ergeben. Sind $\alpha+\beta i,\alpha'+\beta'i,\alpha''+\beta''i$ usw. die bei Division kongruenter Zahlen durch den Modul entstehenden Quotienten, so sind offenbar sowohl die Differenzen der Größen $\alpha,\alpha',\alpha''$ usw. untereinander als auch die Differenzen zwischen den Größen $\beta,\beta',\beta''$ usw. ganze Zahlen. Es ist klar, dass stets ein Rest vorhanden ist, für den $\alpha$ und $\beta$ zwischen den Grenzen $0$ und $1$ liegen, die erstere eingeschlossen, die letztere ausgeschlossen. Einen solchen Rest nennen wir einfach einen kleinsten Rest. Wenn man es vorzieht, kann man stattdessen auch die Grenzen $-\frac12$ und $+\frac12$ wählen, die eine eingeschlossen, die andere ausgeschlossen; einen dieser Beschränkung entsprechenden Rest nennen wir \emph{absolut kleinsten}.

Bezüglich\textperiodcentered dieser kleinsten Reste ergeben sich folgende Aufgaben.

\Article{43}
Den kleinsten Rest einer gegebenen komplexen Zahl $A+Bi$ nach dem Modul $a+bi$, mit der Norm $=p$, findet man wie folgt. Ist $x+yi$ der gesuchte kleinste Rest, so ist $(x+yi)(a-bi)$ der kleinste Rest des Produkts $(A+Bi)(a-bi)$ nach dem Modul $(a+bi)(a-bi)$, also nach dem Modul $p$. Wir setzen daher
\[
aA+bB=Fp+f,\quad aB-bA=Gp+g
\]
sodass $f,g$ die kleinsten Reste der Zahlen $aA+bB,aB-bA$ nach dem Modul $p$ sind. Dann ist
\SourceLeaf{124}%
\[
x+yi=\frac{f+gi}{a-bi}
\]
oder
\begingroup\postdisplaypenalty=10000\relax
\[
\begin{aligned}
x&=\frac{af-bg}{p}=A-aF+bG\\
y&=\frac{ag+bf}{p}=B-aG-bF
\end{aligned}
\]
Offenbar müssen die kleinsten Reste $f,g$ entweder zwischen den Grenzen $0$ und $p-1$ oder zwischen $-\frac12p$ und $+\frac12p$ genommen werden, je nachdem, ob der einfach kleinste oder der absolut kleinste Rest der komplexen Zahl verlangt wird.\endgroup

\Article{44}
Ein vollständiges System kleinster Reste für einen gegebenen Modul lässt sich auf mehrere Arten aufstellen. Das erste Verfahren bestimmt zunächst die Grenzen, innerhalb deren die reellen Glieder liegen müssen, und gibt dann für jeden Wert innerhalb dieser Grenzen die Grenzen der imaginären Teile an. Das allgemeine Kriterium für einen kleinsten Rest $x+yi$ nach dem Modul $a+bi$ besteht darin, dass sowohl $ax+by=\xi$ als auch $ay-bx=\eta$ zwischen den Grenzen $0$ und $aa+bb$ liegen, wenn es sich um einfach kleinste Reste handelt, oder zwischen $-\frac12(aa+bb)$ und $+\frac12(aa+bb)$, wenn absolut kleinste Reste verlangt werden; dabei ist eine Grenze ausgeschlossen. Besondere Regeln würden die Unterscheidung der Fälle erfordern, die aus den verschiedenen Vorzeichen der Zahlen $a,b$ entstehen. Da die Entwicklung dieser Unterscheidung keine Schwierigkeit bietet, halten wir uns hier nicht dabei auf; es genüge, die Art des Verfahrens an einem einzigen Beispiel darzulegen.

Für den Modul $5+2i$ müssen die einfach kleinsten Reste $x+yi$ so beschaffen sein, dass sowohl $5x+2y=\xi$ als auch $5y-2x=\eta$ einer der Zahlen $0,1,2,3\dotsfour28$ gleich ist. Die Gleichung $29x=5\xi-2\eta$ zeigt, dass die positiven Werte von $x$ nicht größer als $\frac{5\pdot28}{29}$ und die negativen, vom Vorzeichen abgesehen, nicht größer als $\frac{2\pdot28}{29}$ sein können. Alle zulässigen Werte von $x$ sind daher $-1,0,1,2,3,4$. Für $x=-1$ muss $2y$ einer der Zahlen $5,6,7\dotsfour33$ und $5y$ einer der Zahlen $-2,-1,0,1\dotsfour26$ gleich sein; somit ist der kleinste Wert von $y$ $+3$, der größte $+5$. Behandelt man die übrigen Werte von $x$ ebenso, so entsteht die folgende Tabelle aller kleinsten Reste:

\SourceLeaf{125}%
\[
\begin{array}{r|l}
x&y\\\hline
-1&3,4,5\\
0&0,1,2,3,4,5\\
+1&1,2,3\pdot4,5\pdot6\\
+2&1,2,3,4,5,6\\
+3&2,3,4\pdot5,6\\
+4&2,3,4
\end{array}
\]

\Needspace{20\baselineskip}
Ebenso müssen für die absolut kleinsten Reste $\xi$ und $\eta$ einer der Zahlen $-14,-13,-12\dotsfour+14$ gleich sein. Daher kann $29x$ nicht außerhalb der Grenzen $-7\pdot14$ und $+7\pdot14$ liegen, und somit muss $x$ einer der Zahlen $-3,-2,-1,0,1,2,3$ gleich sein. Für $x=-3$ ist $2y=\xi-5x=\xi+15$ einer der Zahlen $1,2,3\dotsfour29$ gleich, $5y=\eta+2x=\eta-6$ dagegen einer der Zahlen $-20,-19,-18\dotsfour+8$; daraus ergibt sich für $y$ der einzige Wert $+1$. Behandeln wir die übrigen Werte von $x$ auf dieselbe Weise, so erhalten wir die Tabelle aller absolut kleinsten Reste:

\[
\begin{array}{r|l}
x&y\\\hline
-3&+1\\
-2&-2,-1,0,+1,+2\\
-1&-3,-2,-1,0,+1,+2\\
0&-2,-1,0,+1,+2\\
+1&-2,-1,0,+1,+2,+3\\
+2&-2,-1,0,+1,+2\\
+3&-1
\end{array}
\]

